% Generated draft: requires mathematical and source review.
\documentclass[11pt]{article}
\usepackage[margin=25mm]{geometry}
\usepackage{fontspec}
\setmainfont{Latin Modern Roman}
\usepackage{amsmath,amssymb,mathtools,unicode-math}
\setmathfont{Latin Modern Math}
\usepackage{enumitem,needspace,xurl,hyperref}
\setlength{\parindent}{0pt}
\setlength{\parskip}{5pt}
\setlength{\emergencystretch}{3em}
\newcommand{\N}{\mathbb{N}}
\newcommand{\Z}{\mathbb{Z}}
\newcommand{\Q}{\mathbb{Q}}
\newcommand{\R}{\mathbb{R}}
\newcommand{\C}{\mathbb{C}}
\newcommand{\F}{\mathbb{F}}
\newcommand{\A}{\mathbb{A}}
\newcommand{\PP}{\mathbb{P}}
\newcommand{\E}{\mathbb{E}}
\newcommand{\eps}{\varepsilon}
\newcommand{\dd}{\,\mathrm{d}}
\begin{document}
\section*{UnsolvedMath ID: 800176}
\textbf{Source problem code:} catalogue entry
\subsection{Definitions and mathematical statement}
Let \(\zeta(s)\) be the meromorphic continuation of the Riemann zeta function \(\sum_{n=1}^{\infty}n^{-s}\), initially defined for \(\operatorname{Re}s>1\). The Riemann Hypothesis asserts that every nontrivial zero of \(\zeta\) has real part \(1/2\). Use natural logarithms below.

The supplied source attributes to Chandee and Soundararajan [CS] the following conditional estimate: assuming the Riemann Hypothesis,
\[
\left|\zeta\!\left(\tfrac12+it\right)\right|
\ll \exp\!\left(\left(C_0+o(1)\right)\frac{\log t}{\log\log t}\right)
\qquad(t\to+\infty),
\qquad C_0=\frac{\log 2}{2}.
\]
Here \(t\) is real, \(o(1)\) tends to zero as \(t\to+\infty\), and \(\ll\) denotes an inequality up to a positive multiplicative constant independent of \(t\).

The requested improvement, using the source's exact symbolic value of \(C_0\), is to prove that there exists a real constant
\[
C<\frac{\log 2}{2}
\]
such that, assuming the Riemann Hypothesis, for every \(\varepsilon>0\) there exist constants \(A_{\varepsilon}>0\) and \(T_{\varepsilon}>e\) for which every real \(t\ge T_{\varepsilon}\) satisfies
\[
\left|\zeta\!\left(\tfrac12+it\right)\right|
\le A_{\varepsilon}\exp\!\left((C+\varepsilon)\frac{\log t}{\log\log t}\right).
\]
The constants may depend on the chosen \(C\) and on \(\varepsilon\), but not on \(t\).

\textbf{Unresolved source discrepancy.} The source also writes \(C_0=0.1505\ldots\), whereas with natural logarithms \((\log 2)/2=0.346573\ldots\). These values are incompatible. The statement above provisionally follows the exact symbolic threshold; the intended numerical threshold requires source clarification.
\subsection{Short English statement}
Assuming the Riemann Hypothesis, improve the leading constant in the stated upper bound for \(\lvert\zeta(1/2+it)\rvert\) as real \(t\to+\infty\). The source's symbolic constant \((\log 2)/2\) conflicts with its decimal approximation \(0.1505\ldots\); this draft provisionally uses the symbolic value and requires review.
\subsection{Sources}
\begin{itemize}
\item Supplied local UnsolvedMath registry record: permanent ID \(800176\), exact source problem code \texttt{catalogue entry}; title: ``Moments of L-functions --- Moments of the Riemann zeta function --- Problem 27.'' Upstream revision: \texttt{d371ee917ac9647fbb06cd6fc529deec6f3dda82}.
\item Hamieh, Kadiri, Ng, \emph{Moments of L-functions Problem List} (2024), Problem 27: \url{https://arxiv.org/abs/2405.17800}. The supplied provenance records access on 2026-10-16 and extraction by \texttt{arxiv\_latex}; these are supplied provenance annotations, not a claim of independent source access or verification.
\item The registry lists Hamieh, Kadiri, Ng as proposers, with proposed year 2024. The extracted problem itself carries the attribution ``N. Ng, Lethbridge.''
\item The extracted statement attributes the existing conditional estimate to Chandee and Soundararajan, reference label [CS]. No full bibliographic entry or separate link for [CS] is supplied.
\item The local record labels the problem \texttt{open}. This is a supplied registry classification, not an independently verified assessment of the literature.
\end{itemize}
\end{document}
